\documentclass{article}
\usepackage[T1]{fontenc}
\usepackage[utf8]{inputenc}
\usepackage{lmodern}
\usepackage[a4paper, margin=1in]{geometry}
\usepackage[portuguese]{babel}
\usepackage{url}

\usepackage{minted}
\large
\title{Implementação de Grafos em C++}
\begin{document}
\begin{titlepage}
	\begin{center}
    \line(1,0){300}\\
    [0.65cm]
    \huge{\bfseries Implementação de Grafos em C++}\\
	\line(1,0){300}\\
    \textsc{\LARGE Universidade Estadual do Oeste do Paraná}\\
	\textsc{\Large 3ª Avaliação - Estrutura de Dados}\\
    \textsc{\Large Professor MSc. Leonardo Medeiros}\\
	\textsc{\LARGE \today}\\
    [5.5cm]     
	\end{center}
	\begin{flushright}
		\textsc{\Large Vinícius Almeida Hermes\\vinicius.hermes@unioeste.br}\\
	\end{flushright}
\end{titlepage}

\section*{Introdução}

    Este trabalho consiste na documentação de duas implementações da estrutura de grafos, para obtenção da nota pertinente à terceira avaliação da disciplina de Estrutura de Dados, ministrada pelo professor MSc. Leonardo Medeiros. Com isso, estão documentadas as implementações de grafos por matriz de adjacência e por lista de adjacência, ambas utilizando a linguagem C++.

\section*{Grafo por Matriz de Adjacência}
    Uma das formas de representar um grafo é utilizando uma matriz de adjacência. Dado um grafo G com $n$ vértices, podemos representá-lo em uma matriz \textit{n x n} $A(G)=[a_{ij}]$. Cada índice $a_{ij}$ da matriz armazena o peso da aresta que conecta $a_i$ a $a_j$. Se $a_{ij} = 0$, não há aresta ligando estes dois vértices \cite{wikipedia:matriz_adjacencia}.

\subsection*{Implementação}
\subsubsection*{Definição da classe}

    \textbf{Descrição:} A classe armazena o peso de cada aresta em uma matriz de tamanho máximo de 100 x 100. Cada linha representa o vértice de origem e cada coluna representa o vértice de destino. Uma posição com valor zero indica que não existe aresta entre os respectivos vértices. Já uma posição com valor diferente de zero representa o peso da aresta. Os vértices são identificados por índices no intervalo [0, numero\_vertices - 1].

\textbf{Arquivo:} \texttt{grafo\_matriz.hpp}
\begin{minted}[frame=lines, linenos, fontsize=\small]{cpp}
class Grafo
{
private:
    TipoPeso Matriz[MAX_VERTICES][MAX_VERTICES];
    int N_Vertices;
    int N_Arestas;

public:
    explicit Grafo(int numero_vertices);

    ~Grafo() = default;

    void insereAresta(TipoValorVertice*v1, TipoValorVertice *v2,
                        TipoPeso *peso); 

    bool existeAresta(TipoValorVertice v1, TipoValorVertice v2);

    bool listaAdjVazia(TipoValorVertice *v);

    TipoApontador primeiroListaAdj(TipoValorVertice *v);
    
    void proxAdj(TipoValorVertice *v, TipoValorVertice *adj, TipoPeso *p,
                        TipoApontador *prox, bool *fimListaAdj);
    
    void retiraAresta( TipoValorVertice *v1, TipoValorVertice *v2, TipoPeso *p);

    void imprimeGrafo();
};
\end{minted}

\subsubsection*{Tipos auxiliares e constantes}

\begin{itemize}
    \item \texttt{MAX\_VERTICES}: \textit{limita o tamanho físico máximo para a matriz de adjacência}.
    \item \texttt{TipoValorVertice}: \textit{abstrai o tipo int para definir valores de vértice}.
    \item \texttt{TipoPeso}: \textit{abstrai o tipo int para definir o peso de uma aresta}.
    \item \texttt{TipoApontador}: \textit{abstrai o tipo int para apontar valores para um vértice}.
\end{itemize}

\subsubsection*{Atributos privados}

    Os atributos armazenam o estado interno do objeto.

\begin{itemize}
    \item \texttt{Matriz}: \textit{armazena as arestas e seus pesos da matriz de adjacência}.
    \item \texttt{N\_Vertices}: \textit{armazena o número de vértices}.
    \item \texttt{N\_Arestas}: \textit{armazena o número de arestas}.
\end{itemize}

\subsubsection*{Métodos públicos}
\paragraph{Construtor \texttt{Grafo()}}

    \mbox{}\par
    \noindent
    \textbf{Finalidade:} Recebe o número de vértices e inicializa um grafo com nenhuma aresta.

    \textbf{Parâmetros:}
    \begin{itemize}
        \item \texttt{numero\_vertices}: \textit{número de vértices que se deseja inserir no grafo}.
    \end{itemize}

    \textbf{Retorno:} construtores não possuem valor de retorno.

    \textbf{Funcionamento:} Utiliza dois laços para percorrer o atributo Matriz, atribuindo 0 a cada índice Matriz[i][j] (nenhum peso de aresta).

\begin{minted}[frame=lines, linenos, fontsize=\small]{cpp}
Grafo::Grafo(int numero_vertices)
    : N_Vertices(numero_vertices), N_Arestas(0){
    for (int i = 0; i < numero_vertices; i++){
        for (int j = 0; j < numero_vertices; j++){
            Matriz[i][j] = 0;
        }
    }
}
\end{minted}
\paragraph{Método \texttt{insereAresta()}}

    \mbox{}\par
    \noindent
    \textbf{Finalidade:} Inserir uma aresta (peso) entre dois vértices.

    \textbf{Parâmetros:}
    \begin{itemize}
        \item \texttt{v1}: \textit{vértice de origem}.
        \item \texttt{v2}: \textit{vértice de destino}.
        \item \texttt{peso}: \textit{peso da aresta}.
    \end{itemize}

    \textbf{Retorno:} A função não possui retorno.

    \textbf{Funcionamento:} A função valida os parâmetros, rejeita pesos menores ou iguais a zero e atribui o peso informado à posição correspondente da Matriz. Caso a posição ainda não contenha uma aresta, incrementa o contador de arestas.

\begin{minted}[frame=lines, linenos, fontsize=\small]{cpp}
void Grafo::insereAresta(TipoValorVertice *v1, TipoValorVertice *v2, TipoPeso *peso){
    if (v1 == nullptr || v2 == nullptr || peso == nullptr ||
        *v1 < 0 || *v1 >= N_Vertices ||
        *v2 < 0 || *v2 >= N_Vertices || *peso <= 0){
        return;
    }
    if (this->Matriz[*v1][*v2] == 0){
        ++N_Arestas;
    }
    this->Matriz[*v1][*v2] = *peso;
}
\end{minted}
\paragraph{Método \texttt{existeAresta()}}

    \mbox{}\par
    \noindent
    \textbf{Finalidade:} Verifica se existe uma aresta entre dois pontos quaisquer.

    \textbf{Parâmetros:}
    \begin{itemize}
        \item \texttt{v1}: \textit{vértice de origem}.
        \item \texttt{v2}: \textit{vértice de destino}.
    \end{itemize}

    \textbf{Retorno:} a função retorna \textit{true} caso exista uma aresta entre os vértices (peso > 0). Caso contrário, retorna \textit{false}. O grafo considera apenas pesos positivos.

    \textbf{Funcionamento:} A função acessa a Matriz no índice de parâmetro e verifica se o peso é maior que 0.

\begin{minted}[frame=lines, linenos, fontsize=\small]{cpp}
bool Grafo::existeAresta(TipoValorVertice v1, TipoValorVertice v2){
    if (this->Matriz[v1][v2] > 0){
        return true;
    }
    return false;
}
\end{minted}
\paragraph{Método \texttt{listaAdjVazia()}}

    \mbox{}\par
    \noindent
    \textbf{Finalidade:} Verifica se há alguma adjacência no vértice em questão.

    \textbf{Parâmetros:}
    \begin{itemize}
        \item \texttt{v}: \textit{vértice a ser verificado}.
    \end{itemize}

    \textbf{Retorno:} a função retorna \textit{true} caso não haja nenhuma adjacência. Caso contrário, retorna \textit{false}.

    \textbf{Funcionamento:} A função utiliza um laço para percorrer todas as colunas da linha informada no parâmetro.

\begin{minted}[frame=lines, linenos, fontsize=\small]{cpp}
bool Grafo::listaAdjVazia(TipoValorVertice *v){
    bool vazio = true;
    int j = 0;
    while (j < this->N_Vertices && vazio){
        if (this->Matriz[*v][j] > 0){
            return false;
        } else{
            j++;
        }
    }
    return true;
}
\end{minted}
\paragraph{Método \texttt{primeiroListaAdj()}}

    \mbox{}\par
    \noindent
    \textbf{Finalidade:} Obtém o primeiro vértice que tem adjacência a determinado vértice.

    \textbf{Parâmetros:}
    \begin{itemize}
        \item \texttt{v}: \textit{vértice a ser verificado}.
    \end{itemize}

    \textbf{Retorno:} Índice do primeiro vértice adjacente. Quando não há adjacência, retorna o número de vértices, pois percorreu toda a matriz.

    \textbf{Funcionamento:} A função apresenta exatamente o mesmo comportamento do método \textbf{\textit{listaAdjVazia()}}, mas, ao final, retorna o índice usado para percorrer as colunas. Caso percorra toda a matriz, também exibe uma mensagem de que a lista de adjacência para este vértice está vazia.

\begin{minted}[frame=lines, linenos, fontsize=\small]{cpp}
TipoApontador Grafo::primeiroListaAdj(TipoValorVertice *v){
    TipoApontador j = 0;
    bool vazio = true;

    while(j < this->N_Vertices && vazio){
        if (this->Matriz[*v][j] > 0){
            vazio = false;
        } else{
            j++;
        }
    }

    if(j == this->N_Vertices){
        printf("ERRO: Lista adjacência vazia!\n");
    }

    return j;
}
\end{minted}
\paragraph{Método \texttt{proxAdj()}}

    \mbox{}\par
    \noindent
    \textbf{Finalidade:} Obtém uma adjacência e avança o apontador da lista.

    \textbf{Parâmetros:}
    \begin{itemize}
       \item \texttt{v}: \textit{vértice cuja lista de adjacência será percorrida.}

        \item \texttt{adj}\textsubscript{out}: \textit{recebe o vértice adjacente encontrado.}

        \item \texttt{p}\textsubscript{out}: \textit{Recebe o peso da aresta entre `v` e `adj`.}

        \item \texttt{prox}\textsubscript{in,out}: \textit{Recebe a posição da próxima adjacência a ser consultada.}

        \item \texttt{fimListaAdj}\textsubscript{out}: \textit{Recebe true quando não existem mais adjacências a consultar e false caso ainda existam.}
        
    \end{itemize}

    \textbf{Retorno:} A função não possui retorno. Alguns valores de parâmetros são modificados durante sua execução.

    \textbf{Funcionamento:} A função obtém a adjacência indicada pelo apontador \textit{prox}, armazenando em \textit{adj} o vértice adjacente e em \textit{p} o peso da aresta correspondente. Em seguida, avança até a próxima posição com uma aresta e informa, por meio de \textit{fimListaAdj}, se o final da lista de adjacências foi alcançado.

\begin{minted}[frame=lines, linenos, fontsize=\small]{cpp}
void Grafo::proxAdj(TipoValorVertice *v, TipoValorVertice *adj, TipoPeso *p,
                    TipoApontador *prox, bool *fimListaAdj){
    *adj = *prox;
    *p = this->Matriz[*v][*prox];
    (*prox)++;

    while(*prox < this->N_Vertices && this->Matriz[*v][*prox] == 0){
        (*prox)++;
    }

    if (*prox == this->N_Vertices){
        *fimListaAdj = true;
    } else{
        *fimListaAdj = false;
    }
}
\end{minted}
\paragraph{Método \texttt{retiraAresta()}}

    \mbox{}\par
    \noindent
    \textbf{Finalidade:} Remove uma aresta entre dois vértices e retorna o peso que tinha.

    \textbf{Parâmetros:}
    \begin{itemize}
       \item \texttt{v1}: \textit{vértice de origem.}
       \item \texttt{v2}: \textit{vértice de destino.}

        \item \texttt{p}\textsubscript{out}: \textit{recebe o peso da aresta removida.}
        
    \end{itemize}

    \textbf{Retorno:} A função não possui retorno. Um valor de parâmetro é modificado durante sua execução.

    \textbf{Funcionamento:} A função acessa a matriz no índice Matriz[v1][v2], atribui o valor do peso desta aresta a \textit{p} e em seguida atribui o valor 0 a este índice.

\begin{minted}[frame=lines, linenos, fontsize=\small]{cpp}
void Grafo::retiraAresta( TipoValorVertice *v1, TipoValorVertice *v2, TipoPeso *p){
    if (v1 == nullptr || v2 == nullptr || p == nullptr ||
        *v1 < 0 || *v1 >= N_Vertices ||
        *v2 < 0 || *v2 >= N_Vertices ||
        this->Matriz[*v1][*v2] == 0){
        printf("Aresta não existe!\n");
    } else{
        *p = this->Matriz[*v1][*v2];
        this->Matriz[*v1][*v2] = 0;
        --N_Arestas;
    }
}
\end{minted}






\paragraph{Método \texttt{imprimeGrafo()}}

    \mbox{}\par
    \noindent
    \textbf{Finalidade:} Imprime a matriz de adjacência do grafo, com os pesos de cada aresta.

        \textbf{Parâmetros:} A função não possui parâmetros.

    \textbf{Retorno:} A função não possui retorno.

    \textbf{Funcionamento:} A função imprime primeiro os índices das colunas. Em seguida, percorre com dois laços a matriz de adjacência, antes de começar a imprimir cada peso na linha, o laço imprime o índice atual, para formatar a tabela da matriz de adjacência.

\begin{minted}[frame=lines, linenos, fontsize=\small]{cpp}
void Grafo::imprimeGrafo(){
    printf("   ");
    for (int i = 0; i < this->N_Vertices; i++){
        printf("%3d", i);
    }
    printf("\n");

    for (int i = 0; i < this->N_Vertices; i++){
        printf("%3d", i);
        for (int j = 0; j < this->N_Vertices; j++){
            printf("%3d", this->Matriz[i][j]);
        }
        printf("\n");
    }
    cout << endl;
}
\end{minted}






























\section*{Grafo por Lista de Adjacência}
Em uma representação de grafos utilizando lista de adjacência, podemos manter, para cada vértice do grafo, uma lista de todos os outros vértices com os quais ele tem uma aresta \cite{wikipedia:lista_adjacencia}.

\subsection*{Implementação}

\subsubsection*{Definição de classes auxiliares e seus métodos}

    \textbf{Descrição:} Para o funcionamento do grafo por lista de adjacência, foi necessária a criação de classes auxiliares. O grafo consiste em uma "lista de listas", e cada uma dessas listas é composta por células que, cada uma delas armazena um item (que armazena o peso e o vértice) e o ponteiro para a próxima célula.
\textbf{Arquivo:} \texttt{grafo\_lista.hpp}
\begin{minted}[frame=lines, linenos, fontsize=\small]{cpp}

class TipoItem
{
public:
    TipoValorVertice Vertice;
    TipoPeso Peso;
    // Definição de seu construtor
    TipoItem(TipoValorVertice v = 0, TipoPeso p = 0): Vertice(v), Peso(p){}
};


class TipoCelula {
public:
    TipoItem item;
    TipoCelula* prox;
    // Definição de seu construtor
    TipoCelula(const TipoItem& item = TipoItem(), TipoCelula* prox = nullptr)
        : item(item), prox(prox){}
};

class TipoLista{
public:
    TipoApontador Primeiro;
    TipoApontador Ultimo;
    TipoLista();
    ~TipoLista();
    bool vazia() const;
    void insere(const TipoItem& i);
    void insere(const TipoItem& i, TipoApontador l);
    bool retira(TipoApontador p, TipoItem *i);
    void imprime() const;
};
\end{minted}

\paragraph*{Tipos auxiliares e constantes}

\begin{itemize}
    \item \texttt{MAX\_VERTICES}: \textit{limita o tamanho físico máximo para as listas de adjacência}.
    \item \texttt{TipoValorVertice}: \textit{abstrai o tipo int para definir valores de vértice}.
    \item \texttt{TipoPeso}: \textit{abstrai o tipo int para definir o peso de uma aresta}.
    \item \texttt{TipoApontador}: \textit{abstrai o ponteiro utilizado para apontar para uma célula}.

    
\end{itemize}

\paragraph*{Métodos da classe TipoLista}
\subparagraph{Construtor \texttt{TipoLista()}}

    \mbox{}\par
    \noindent
    \textbf{Finalidade:} Constrói uma lista vazia.

    \textbf{Parâmetros:} A função não possui parâmetros.

    \textbf{Retorno:} construtores não possuem valor de retorno.

    \textbf{Funcionamento:} Apenas utiliza o operador \textit{new} para inicializar a lista.

\begin{minted}[frame=lines, linenos, fontsize=\small]{cpp}
TipoLista::TipoLista()
    : Primeiro(new TipoCelula()), Ultimo(Primeiro){}
\end{minted}

\subparagraph{Método \texttt{FLVazia()}}

    \mbox{}\par
    \noindent
    \textbf{Finalidade:} Torna a lista vazia.

    \textbf{Parâmetros:} A função não possui parâmetros.

    \textbf{Retorno:} A função não possui retorno.

    \textbf{Funcionamento:} Percorre toda a lista, deletando a célula atual a cada iteração. 

\begin{minted}[frame=lines, linenos, fontsize=\small]{cpp}
void TipoLista::FLVazia(){
    TipoApontador atual = Primeiro->prox;
    while (atual != nullptr){
        TipoApontador proxima = atual->prox;
        delete atual;
        atual = proxima;
    }
    Primeiro->prox = nullptr;
    Ultimo = Primeiro;
}
\end{minted}

\subparagraph{Método \texttt{vazia()}}

    \mbox{}\par
    \noindent
    \textbf{Finalidade:} Verifica se a lista está vazia.

    \textbf{Parâmetros:} A função não possui parâmetros.

    \textbf{Retorno:} a função retorna \textit{true} caso a lista esteja vazia. Caso contrário, retorna \textit{false}.

    \textbf{Funcionamento:} A função apenas verifica se a última célula é igual a primeira, pois se a afirmação é verdadeira é porque não há elementos.

\begin{minted}[frame=lines, linenos, fontsize=\small]{cpp}
bool TipoLista::vazia() const {
    return Primeiro == Ultimo;
}
\end{minted}



\subparagraph{Método \texttt{insere()}}

    \mbox{}\par
    \noindent
    \textbf{Finalidade:} Insere elementos na lista.

    \textbf{Nota:} Utiliza de sobrecarga para definir se irá inserir no final da lista ou em uma posição específica.

    \textbf{Parâmetros:}
    \begin{itemize}
       \item \texttt{i}: \textit{item a ser inserido.}
       \item \texttt{l}: \textit{posição a ser inserida (opcional - utilizada apenas quando deseja-se escolher a posição da inserção).}        
    \end{itemize}

    \textbf{Retorno:} A função não possui retorno.

    \textbf{Funcionamento:} A função utiliza o apontador como índice para criar uma nova célula na posição, apontando a anterior (posição) para esta nova e apontando a nova para a próxima, que era da anterior. Em caso de a posição ser a última, aponta o valor \texttt{Ultimo} da lista para a nova célula.

\begin{minted}[frame=lines, linenos, fontsize=\small]{cpp}
void TipoLista::insere(const TipoItem& i){
    insere(i, Ultimo);
}

void TipoLista::insere(const TipoItem& i, TipoApontador l){
    if (l == nullptr){
        return;
    }
    TipoApontador nova = new TipoCelula(i, l->prox);
    l->prox = nova;
    if (l == Ultimo){
        Ultimo = nova;
    }
}
\end{minted}



\subparagraph{Método \texttt{retira()}}

    \mbox{}\par
    \noindent
    \textbf{Finalidade:} Retira elementos da lista.

    \textbf{Parâmetros:}
    \begin{itemize}
       \item \texttt{p}: \textit{elemento anterior ao que será removido.}
       \item \texttt{i}: \textit{receberá durante a execução o item que foi removido, contendo seu peso e seu vértice.}        
    \end{itemize}

    \textbf{Retorno:} a função retorna \textit{true} caso a remoção seja bem-sucedida. Caso contrário, retorna \textit{false}.

    \textbf{Funcionamento:} A função cria um auxiliar que recebe o próximo de \textit{p}, altera \textit{i} para receber os atributos do item que será removido, e deleta-o. A função também verifica caso o item a ser removido seja o último da lista, para alterar o apontador de último.

\begin{minted}[frame=lines, linenos, fontsize=\small]{cpp}
bool TipoLista::retira(TipoApontador p, TipoItem *i){
    if (p == nullptr || p->prox == nullptr || i == nullptr){
        return false;
    }
    TipoApontador q = p->prox;
    *i = q->item;
    p->prox = q->prox;
    if (q == Ultimo){
        Ultimo = p;
    }
    delete q;
    return true;
}
\end{minted}





\subparagraph{Método \texttt{imprime()}}

    \mbox{}\par
    \noindent
    \textbf{Finalidade:} Imprime a lista.

    \textbf{Parâmetros:} A função não possui parâmetros.

    \textbf{Retorno:} A função não possui retorno.

    \textbf{Funcionamento:} Percorre a lista imprimindo cada vértice e peso.

\begin{minted}[frame=lines, linenos, fontsize=\small]{cpp}
void TipoLista::imprime() const {
    TipoApontador aux = Primeiro->prox;
    while (aux != nullptr){
        printf("%3d (%d) ", aux->item.Vertice, aux->item.Peso);
        aux = aux->prox;
    }
    putchar('\n');
}
\end{minted}










%%%%%%%%%%%%%%%%%%%%%%%%%%%%%

\subsubsection*{Definição de classe TipoGrafo}
    \textbf{Descrição:} Define o grafo, contendo as listas de células, o número de vértices e arestas.
\textbf{Arquivo:} \texttt{grafo\_lista.hpp}
\begin{minted}[frame=lines, linenos, fontsize=\small]{cpp}
class TipoGrafo
{
public:
    TipoLista Adj[MAX_VERTICES + 1];
    TipoValorVertice N_Vertices;
    short N_Arestas;

    TipoGrafo(TipoValorVertice n);

    void FGVazio();

    void insereAresta(TipoValorVertice *v1, TipoValorVertice *v2, TipoPeso *p);

    bool existeAresta(TipoValorVertice v1, TipoValorVertice v2);

    bool listaAdjVazia(TipoValorVertice *v);

    TipoApontador primeiroListaAdj(TipoValorVertice *v);

    void proxAdj(TipoValorVertice *v, TipoValorVertice *adj, TipoPeso *p, TipoApontador                     *prox, bool *fimListaAdj);

    void retiraAresta(TipoValorVertice *v1, TipoValorVertice *v2, TipoPeso *p);

    void liberaGrafo();

    void imprimeGrafo();
};
\end{minted}

\paragraph*{Atributos da classe}

    Os atributos armazenam o estado interno do objeto.

\begin{itemize}
    \item \texttt{Adj}: \textit{vetor que armazena as listas de adjacência}.
    \item \texttt{N\_Vertices}: \textit{armazena o número de vértices}.
    \item \texttt{N\_Arestas}: \textit{armazena o número de arestas}.
\end{itemize}


\paragraph*{Métodos da classe TipoGrafo}

\subparagraph{Método \texttt{TipoGrafo()}}

    \mbox{}\par
    \noindent
    \textbf{Finalidade:} Constrói objeto TipoGrafo.

    \textbf{Parâmetros:}
    \begin{itemize}
       \item \texttt{n}: \textit{número de vértices que o grafo terá.}
    \end{itemize}

    \textbf{Retorno:} construtores não possuem valor de retorno.

    \textbf{Funcionamento:} Atribui o número de vértices do parâmetro ao N\_Vertices do objeto.

\begin{minted}[frame=lines, linenos, fontsize=\small]{cpp}
TipoGrafo::TipoGrafo(TipoValorVertice n)
    : N_Vertices(n), N_Arestas(0) {
    if (n < 0 || n > MAX_VERTICES) {
        N_Vertices = 0;
    }
}
\end{minted}











\subparagraph{Método \texttt{FGVazio()}}

    \mbox{}\par
    \noindent
    \textbf{Finalidade:} Torna o grafo vazio.

    \textbf{Parâmetros:} A função não possui parâmetros.

    \textbf{Retorno:} A função não possui retorno.

    \textbf{Funcionamento:} Percorre o vetor de listas de adjacência, e limpa cada uma delas.

\begin{minted}[frame=lines, linenos, fontsize=\small]{cpp}
void TipoGrafo::FGVazio(){
    for (int i = 0; i < N_Vertices; i++){
        Adj[i].FLVazia();
    }
    N_Arestas = 0;
}
\end{minted}


\subparagraph{Método \texttt{insereAresta()}}

    \mbox{}\par
    \noindent
    \textbf{Finalidade:} Insere uma aresta no grafo.

    \textbf{Parâmetros:}
    \begin{itemize}
        \item \texttt{v1}: \textit{vértice de origem.}
        \item \texttt{v2}: \textit{vértice de destino.}
        \item \texttt{p}: \textit{peso da aresta.}
    \end{itemize}

    \textbf{Retorno:} A função não possui retorno.

    \textbf{Funcionamento:} Insere na lista do vértice seu destino com seu peso, através do método \textit{insere()} do TipoLista.

\begin{minted}[frame=lines, linenos, fontsize=\small]{cpp}
void TipoGrafo::insereAresta(TipoValorVertice *v1, TipoValorVertice *v2, TipoPeso *p) {
    if (v1 == nullptr || v2 == nullptr || p == nullptr ||
        *v1 < 0 || *v1 >= N_Vertices ||
        *v2 < 0 || *v2 >= N_Vertices) {
        return;
    }
    Adj[*v1].insere(TipoItem(*v2, *p));
    ++N_Arestas;
}
\end{minted}

\subparagraph{Método \texttt{existeAresta()}}

    \mbox{}\par
    \noindent
    \textbf{Finalidade:} Verifica se há uma aresta entre dois vértices.

    \textbf{Parâmetros:}
    \begin{itemize}
        \item \texttt{v1}: \textit{vértice de origem.}
        \item \texttt{v2}: \textit{vértice de destino.}
    \end{itemize}

    \textbf{Retorno:} a função retorna \textit{true} caso exista uma aresta entre os vértices. Caso contrário, retorna \textit{false}.

    \textbf{Funcionamento:} Percorre a lista de adjacência do vértice até encontrar o possível vértice de conexão.


\begin{minted}[frame=lines, linenos, fontsize=\small]{cpp}
bool TipoGrafo::existeAresta(TipoValorVertice v1, TipoValorVertice v2){
    if (v1 < 0 || v1 >= N_Vertices || v2 < 0 || v2 >= N_Vertices){
        return false;
    }
    for (TipoApontador aux = Adj[v1].Primeiro->prox;
         aux != nullptr;
         aux = aux->prox) {
        if (aux->item.Vertice == v2){
            return true;
        }
    }
    return false;
}
\end{minted}

\subparagraph{Método \texttt{listaAdjVazia()}}

    \mbox{}\par
    \noindent
    \textbf{Finalidade:} Verifica se a lista de adjacência para o vértice está vazia.

    \textbf{Parâmetros:}
    \begin{itemize}
       \item \texttt{v}: \textit{vértice a ser verificado.}
    \end{itemize}

    \textbf{Retorno:} a função retorna \textit{true} caso não exista nenhuma adjacência para o vértice. Caso contrário, retorna \textit{false}.

    \textbf{Funcionamento:} Chama a função de verificação se a lista está vazia para o índice do vértice.

\begin{minted}[frame=lines, linenos, fontsize=\small]{cpp}
bool TipoGrafo::listaAdjVazia(TipoValorVertice *v){
    if (v == nullptr || *v < 0 || *v >= N_Vertices){
        return true;
    }
    return Adj[*v].vazia();
}
\end{minted}

\subparagraph{Método \texttt{primeiroListaAdj()}}

    \mbox{}\par
    \noindent
    \textbf{Finalidade:} Verifica qual é a primeira aresta da lista de adjacência do vértice.

    \textbf{Parâmetros:}
    \begin{itemize}
       \item \texttt{v}: \textit{vértice a ser verificado.}
    \end{itemize}

    \textbf{Retorno:} o primeiro vértice adjacente ao vértice verificado.

    \textbf{Funcionamento:} Busca o próximo do primeiro vértice da lista (a primeira adjacência).

\begin{minted}[frame=lines, linenos, fontsize=\small]{cpp}
TipoApontador TipoGrafo::primeiroListaAdj(TipoValorVertice *v){
    if (v == nullptr || *v < 0 || *v >= N_Vertices){
        return nullptr;
    }
    return Adj[*v].Primeiro->prox;
}
\end{minted}

\subparagraph{Método \texttt{proxAdj()}}

    \mbox{}\par
    \noindent
    \textbf{Finalidade:} Obtém uma adjacência e avança o apontador da lista.

    \textbf{Parâmetros:}
    \begin{itemize}
       \item \texttt{v}: \textit{vértice cuja lista de adjacência será percorrida.}

        \item \texttt{adj}\textsubscript{out}: \textit{recebe o vértice adjacente encontrado.}

        \item \texttt{p}\textsubscript{out}: \textit{Recebe o peso da aresta entre `v` e `adj`.}

        \item \texttt{prox}\textsubscript{in,out}: \textit{Recebe a posição da próxima adjacência a ser consultada.}

        \item \texttt{fimListaAdj}\textsubscript{out}: \textit{Recebe true quando não existem mais adjacências a consultar e false caso ainda existam.}
        
    \end{itemize}

    \textbf{Retorno:} A função não possui retorno. Alguns valores de parâmetros são modificados durante sua execução.

    \textbf{Funcionamento:} A função obtém a adjacência indicada pelo apontador \textit{prox}, armazenando em \textit{adj} o vértice adjacente e em \textit{p} o peso da aresta correspondente. Em seguida, avança até a próxima posição com uma aresta e informa, por meio de \textit{fimListaAdj}, se o final da lista de adjacências foi alcançado.

\begin{minted}[frame=lines, linenos, fontsize=\small]{cpp}
void TipoGrafo::proxAdj(TipoValorVertice *v, TipoValorVertice *adj, TipoPeso *p,                                        TipoApontador *prox, bool *fimListaAdj){
    if (v == nullptr || adj == nullptr || p == nullptr ||
        prox == nullptr || fimListaAdj == nullptr || *prox == nullptr){
        *fimListaAdj = true;
        return;
    }
    *adj = (*prox)->item.Vertice;
    *p = (*prox)->item.Peso;
    *prox = (*prox)->prox;
    *fimListaAdj = (*prox == nullptr);
}
\end{minted}

\subparagraph{Método \texttt{retiraAresta()}}

    \mbox{}\par
    \noindent
    \textbf{Finalidade:} Remove uma aresta entre dois vértices e coloca em `*p` o peso que a aresta tinha.

    \textbf{Parâmetros:}
    \begin{itemize}
       \item \texttt{v1}: \textit{vértice de origem.}
       \item \texttt{v2}: \textit{vértice de destino.}
       \item \texttt{p}: \textit{receberá o peso da aresta removida.}

    \end{itemize}

    \textbf{Retorno:} A função não possui retorno, mas altera uma variável durante sua execução.

    \textbf{Funcionamento:} Percorre a lista até encontrar os dois vértices que estão ligados pela aresta e remove o item da lista.

\begin{minted}[frame=lines, linenos, fontsize=\small]{cpp}
void TipoGrafo::retiraAresta(TipoValorVertice *v1, TipoValorVertice *v2, TipoPeso *p){
    if (v1 == nullptr || v2 == nullptr || p == nullptr ||
        *v1 < 0 || *v1 >= N_Vertices ||
        *v2 < 0 || *v2 >= N_Vertices){
        return;
    }
    TipoApontador anterior = Adj[*v1].Primeiro;
    while (anterior->prox != nullptr){
        if (anterior->prox->item.Vertice == *v2){
            TipoItem removido;
            if (Adj[*v1].retira(anterior, &removido)){
                *p = removido.Peso;
                --N_Arestas;
            }
            return;
        }
        anterior = anterior->prox;
    }
}
\end{minted}

\subparagraph{Método \texttt{liberaGrafo()}}

    \mbox{}\par
    \noindent
    \textbf{Finalidade:} Limpa o grafo.

    \textbf{Parâmetros:} A função não possui parâmetros.

    \textbf{Retorno:} A função não possui retorno.

    \textbf{Funcionamento:} Chama a função que torna o grafo vazio.

\begin{minted}[frame=lines, linenos, fontsize=\small]{cpp}
void TipoGrafo::liberaGrafo(){
    FGVazio();
}
\end{minted}

\subparagraph{Método \texttt{imprimeGrafo()}}

    \mbox{}\par
    \noindent
    \textbf{Finalidade:} Imprime as listas de adjacência de cada vértice.

    \textbf{Parâmetros:} A função não possui parâmetros.

    \textbf{Retorno:} A função não possui retorno.

    \textbf{Funcionamento:} Para cada vértice, imprime a sua lista de adjacência.

\begin{minted}[frame=lines, linenos, fontsize=\small]{cpp}
void TipoGrafo::imprimeGrafo(){
    for (int i = 0; i < N_Vertices; i++){
        printf("Vertice%2d: ", i);
        Adj[i].imprime();
    }
}
\end{minted}

\section*{Referências}

\begin{thebibliography}{2}

\bibitem{wikipedia:matriz_adjacencia}
MATRIZ DE ADJACÊNCIA. In: WIKIPÉDIA, a enciclopédia livre. Flórida: Wikimedia Foundation, 2017. 
Disponível em: \url{https://wikipedia.org}. Acesso em: 6 jan. 2017.

\bibitem{wikipedia:lista_adjacencia}
LISTA DE ADJACÊNCIA. In: WIKIPÉDIA, a enciclopédia livre. Flórida: Wikimedia Foundation, 2018. 
Disponível em: \url{https://wikipedia.org}. Acesso em: 20 mar. 2018.

\end{thebibliography}

\end{document}
